\textbf{Keywords:}

Processing of signals has shifted into the digital domain. But the real
world is analog. In order to interact with the analog we need to convert
the digital signals (discrete-value, discrete-time) back to analog
signals (continuous value, continuous time).

The SI base units define the fundamental analog quantities as second,
meter, kilogram, ampere, kelvin, mole and candela (see
\href{/aic2026/a_refresher}{the refresher}). Assume that electronic
circuits interact with the real world in terms of second and ampere.

Related to Ampere we have the derived units of charge (Ampere Seconds),
Volt (W/A), Ohm (V/A), or indeed Siemens (1/\(\Omega\)).


{
\centering
\aicfig{media/NIST.SP_.1247.png}
}

\emph{Figure 1: NIST poster of the SI base units and the derived units.
Image: NIST, US Department of Commerce (US federal work)}

As such, to create a digital to analog converter, we somehow have to
create a circuit that has a function of

\[ I_{out} = D_{in} \times I_{ref}\text{ [I]} \]

\[ t_{out} = D_{in} \times t_{ref}\text{ [s]} \]

\[ Q_{out} = D_{in} \times Q_{ref}\text{ [C]} \]

\[ V_{out} = D_{in} \times V_{ref}\text{ [V]} \]

\[ R_{out} = D_{in} \times R_{ref}\text{ [}\Omega\text{]} \]

The digital value is dimensionless, as such, there must be a reference
value

Digital to analog conversion can be indirect through the relations
between voltage, resistance, current, time, inductance and capacitance.

\[ V = R I \]

\[ Q = C V \]

\[ dt = \frac{C dV}{I} \]

\[ dt = \frac{L dI}{V} \]

\section{Resistor based DACs}\label{resistor-based-dacs}\index{resistor based dacs@Resistor based DACs}


{
\centering
\aicfig{media/dac_r_div_tikz.pdf}
}

\emph{Figure 2: 1-bit DAC: three series resistors with transistor
switches selecting the output tap}

\[ I_{ref} = \frac{V_{ref}}{3 R}\]

\[ V_{out} = D_{in} R I_{ref} = \frac{D_{in} R V_{ref}}{3 R} \text{ [V]} \]

\[ V_{out} = \frac{b_0 2 R V_{ref}}{3 R}  + \frac{\overline{b_0} R V_{ref}}{3 R} \]

\[  V_{out} = \begin{cases} 
\frac{2}{3} V_{ref}, & b_0 = 1 \\
\frac{1}{3} V_{ref}, & b_0 = 0 
\end{cases}
\]


{
\centering
\aicfig{media/dac_r_div2_tikz.pdf}
}

\emph{Figure 3: 1-bit DAC with two resistors, selecting either
\(V_{REF}\) or the midpoint}

\[ I_{ref} = \frac{V_{ref}}{2 R}\]

\[ V_{out} = \frac{b_0 2 R V_{ref}}{2 R}  + \frac{\overline{b_0} R V_{ref}}{2 R} \]

\[  V_{out} = \begin{cases} 
 V_{ref}, & b_0 = 1 \\
\frac{1}{2} V_{ref}, & b_0 = 0 
\end{cases}
\]


{
\centering
\aicfig{media/dac_r_div2b_tikz.pdf}
}

\emph{Figure 4: Two possible 2-bit resistor string DACs with switch
trees}

\section{DAC errors}\label{dac-errors}\index{dac errors@DAC errors}

Digital to analog converters do not add quantization error. The
quantization error is already in the digital word.

\[ V_{out} = B + a_1 D_{in} + \left( a_2 D_{in}^2 + \dots + a_n D_{in}^n \right) \]

Read that as three separate defects. \(B\) is offset, a constant added
to every code. \(a_1\) is gain, which stretches the transfer curve but
keeps it straight. Everything in the bracket is non-linearity, and it is
the only part that a calibration of gain and offset cannot remove.

DAC output will contain gain errors, offset errors, and non-linear
components


{
\centering
\aicfig{media/dac_error_tikz.pdf}
}

\emph{Figure 5: Quantization error of a 2-bit DAC. The digital code
(top), two conventions for turning that code back into a voltage
(middle), and the error each one makes (bottom)}


{
\centering
\aicfig{media/dac_inl_dnl_tikz.pdf}
}

\emph{Figure 6: DAC output compared to the ideal straight line, with INL
and DNL versus digital code}

\[ DNL[k] = \frac{V[k+1] - V[k]}{V_{LSB}} - 1 \]

\[ INL[k] = \frac{V[k] - V_{ideal}[k]}{V_{LSB}} \]

\section{DAC complexity}\label{dac-complexity}\index{dac complexity@DAC complexity}


{
\centering
\aicfig{media/dac_r_switches_tikz.pdf}
}

\emph{Figure 7: Binary switch tree for a resistor string DAC}

The resistor string itself is the easy part - \(2^N\) equal resistors
give \(2^N\) perfectly ordered taps, and the DAC is monotonic by
construction. The cost hides in the selection: something must connect
exactly one tap to the output. The obvious structure is a binary tree,
where each bit steers a rank of switches, and the count below says why
nobody stops there.

As number of resistors grow, the switches grow as

\[ \sum_{n=1}^{N} 2^n = 2^{N+1} - 2 \]


{
\centering
\aicfig{media/dac_r_rows_tikz.pdf}
}

\emph{Figure 8: Row and column switch matrix for a resistor string DAC}

The matrix halves the damage by decoding in two dimensions, the way a
memory does: the row decoder picks a group of taps, the column switches
pick one of them. The tap switches are still there - they must be, every
tap has to be reachable - but the tree above them collapses into one
switch per column.

Give every tap a switch onto a column line, and every column line a
switch to the output. For \(2^N\) taps in a square arrangement that is

\[ 2^{N} + 2^{N/2} \]


{
\centering
\aicfig{media/dac_r_segmented_tikz.pdf}
}

\emph{Figure 9: Segmented DAC switch arrangement combining switch
matrices and a tree}

Switches in a 10-bit digital to analog converter.

\[
\begin{array}{lll}
\text{Tree} & 2^{N+1}-2 & = 2046 \\
\text{Matrix} & 2^{N}  + 2^{N/2}  & = 1056 \\
\text{Two strings, 6b + 4b} & 2^{M} + 2^{N-M} & = 80 \\
\text{Two strings, 5b + 5b} & 2^{M} + 2^{N-M} & = 64
\end{array}
\]

Those three lines are not three versions of the same circuit, and the
difference between the first two and the last is the point of this
section.

The tree and the matrix both address one string of \(2^N\) resistors, so
both need at least one switch per tap. The tree adds switches at every
internal node on top of that, which is why it is the worst of the three;
the matrix adds only one per column, which is why it is roughly half the
tree. Neither can go below \(2^N\), because every tap has to be
reachable on its own.

It is worth being clear that combining them does not help. A tree of
matrix blocks --- a plausible reading of the figure above --- still
needs its \(2^N\) tap switches and now pays for the tree as well: a
4-bit tree over sixteen 6-bit matrices comes to 1182, worse than the
plain matrix at 1056. Every split is worse. There is nothing to win by
decoding the same string more cleverly.

The saving comes from not having one string at all. Put a coarse string
of \(2^M\) resistors in series and let a fine string of \(2^{N-M}\)
resistors interpolate between two adjacent coarse taps, and each
selector only has to reach its own string: \(2^M + 2^{N-M}\) switches
rather than \(2^N\). For ten bits split six and four that is 80, and the
best split is five and five at 64 --- against 1056 for the matrix.
Trading one big decoder for two small ones is worth a factor of sixteen
here, and the reason is simply that \(2^M + 2^{N-M}\) grows far more
slowly than \(2^N\).

Nothing is free: the fine string loads the coarse one and disturbs the
very voltage it is interpolating, a real two-string DAC needs a second
switch per coarse tap to bracket the segment, and the matching now has
to hold between two strings rather than within one. But the switch count
is no longer what stops you.

Large number of bits, will be large number of resistors and switches.

\section{Binary scaled DACs}\label{binary-scaled-dacs}\index{binary scaled dacs@Binary scaled DACs}

A string DAC pays \(2^N\) resistors for \(N\) bits. The R-2R ladder pays
\(2N\): each section divides the remaining voltage by two, so the branch
currents come out binary weighted with only two resistor values. The
next four figures build the ladder one property at a time - the
termination, the input resistance that stays 2R at every section, and
the halving branch currents that make it a DAC.

\[ R_{in} = 2R \parallel 2R = R \]


{
\centering
\aicfig{media/dac_2r_0_tikz.pdf}
}

\emph{Figure 10: R-2R ladder termination: two 2R resistors in parallel
equal R}

\[ R_{in} = R + R = 2R \]

\[ I_{0} = \frac{V_0}{2R} = \frac{V_1}{4R} \]


{
\centering
\aicfig{media/dac_2r_1_tikz.pdf}
}

\emph{Figure 11: One R-2R ladder section: the series R makes the input
resistance 2R}

\[ R_{in} = 2R \parallel 2R = R \]

\[ I_{0} = \frac{V_0}{2R} = \frac{V_1}{4R} \]

\[ I_{1} = \frac{V_1}{2R}\]


{
\centering
\aicfig{media/dac_2r_2_tikz.pdf}
}

\emph{Figure 12: R-2R ladder section with the binary weighted branch
currents \(I_1\) and \(I_0\)}

\[ I_{RF} = I_1b_1 + I_0b_0 = \frac{V_{REF}}{2R}b_1 + \frac{V_{REF}}{4R}b_0 \]

\[ V_{O} = \left(\frac{V_{REF}}{2R}b_1 + \frac{V_{REF}}{4R}b_0\right)R_{F0}\]


{
\centering
\aicfig{media/dac_2r_full_tikz.pdf}
}

\emph{Figure 13: 2-bit R-2R DAC with switched branch currents summed by
a transimpedance amplifier}

The switches steer each branch current either into the virtual ground of
the amplifier or to real ground, so the ladder's currents never change -
only their destination does. That is what makes the R-2R fast for its
size. What it gives up is the string's built-in monotonicity: at the
major transition the MSB branch must match the sum of all the others to
within an LSB, and that is now a matching requirement on the resistors
rather than a property of the structure.

\subsection{Binary coding}\label{binary-coding}\index{binary coding@Binary coding}

For 4 states (2-bit) there are 12 possible transitions


{
\centering
\aicfig{media/dac_bin_states_tikz.pdf}
}

\emph{Figure 14: The 12 possible transitions between the four 2-bit
binary states}

Assume MSB first (left)

\[ 1 \rightarrow 3 \rightarrow 2 \]

Assume LSB first (right)

\[ 1 \rightarrow 0 \rightarrow 2 \]

Both cause a non-monotonic glitch during transition.


{
\centering
\aicfig{media/dac_bin_btran_tikz.pdf}
}

\emph{Figure 15: Binary code transitions with MSB first (left) and LSB
first (right), both non-monotonic}

The switches never move at exactly the same time. Between the old code
and the new one the DAC output visits whatever code the half-switched
bits happen to spell, and around the major transition - 0111 to 1000 -
that intermediate code can be far away. The result is a glitch whose
energy grows with the weight of the bits involved, and no amount of
matching removes it: it is a property of the code, not of the elements.

\subsection{Thermometer encoding}\label{thermometer-encoding}\index{thermometer encoding@Thermometer encoding}


{
\centering
\aicfig{media/dac_thermo_states_tikz.pdf}
}

\emph{Figure 16: Transitions between the thermometer encoded states}

The sequence of MSB to LSB does not matter.

\[ 0 \rightarrow 1 \rightarrow 2  \rightarrow 3\]


{
\centering
\aicfig{media/dac_thermo_tran_tikz.pdf}
}

\emph{Figure 17: Thermometer code transitions are monotonic regardless
of bit order}

Thermometer coding removes the glitch by construction: one more LSB
always means one more element turned on, so the output can only move one
step, whatever order the switches settle in. Monotonicity comes for free
for the same reason. The price is \(2^N - 1\) elements and the decoder
that drives them, which is why real converters segment - thermometer for
the MSBs where the glitch would be worst, binary for the LSBs where it
cannot hurt.


{
\centering
\aicfig{media/dac_r_thermo_tikz.pdf}
}

\emph{Figure 18: Thermometer coded resistor DAC with equal resistors
summed by a transimpedance amplifier}

\section{Current mode DACs}\label{current-mode-dacs}\index{current mode dacs@Current mode DACs}


{
\centering
\aicfig{media/dac_i_tikz.pdf}
}

\emph{Figure 19: Current mode DAC: binary sized differential current
cells switched into a transimpedance output stage}

At high sample rates the resistor structures run out of settling time,
and the current steering DAC takes over: every cell is a current source
that is always on, and the data only chooses which side of a
differential pair the current leaves through. Nothing charges or
discharges except the switch nodes, so this is the architecture behind
every GS/s transmitter DAC.


{
\centering
\aicfig{media/dac_i_vbias_tikz.pdf}
}

\emph{Figure 20: Current mode DAC where the switch drive swings around
\(V_{bias}\) instead of rail to rail}

Driving the steering pair rail to rail briefly turns both switches off
and slams the source node; the cell's current has to go somewhere, and
it goes into the output as a spike. Limiting the switch drive to a small
swing around \(V_{bias}\) keeps the pair in its active region through
the crossover, keeps the current source in saturation, and is the
difference between a DAC that meets its SFDR and one that only meets its
resolution.

\section{Summary}\label{summary}

The one-page version of this chapter:

\begin{itemize}
\tightlist
\item
  A DAC turns a code into charge, current or voltage by summing weighted
  unit elements
\item
  Binary weighting is compact but must switch half the array at the
  major carry; thermometer coding is monotonic and glitch-free but costs
  2\^{}N elements and decoding
\item
  Segmentation spends thermometer coding on the MSBs where it matters
  and binary on the LSBs where it is cheap
\item
  Static accuracy is INL/DNL set by element matching (Pelgrom: area buys
  bits); dynamic accuracy is glitch energy and SFDR
\item
  The references and the switch drivers are part of the DAC: their noise
  and timing skew show up in the output spectrum
\end{itemize}

\section{Would you like to know
more?}\label{would-you-like-to-know-more}

A 28-nm 75-fsrms Analog Fractional-N Sampling PLL With a Highly Linear
DTC Incorporating Background DTC Gain Calibration and Reference Clock
Duty Cycle Correction {[}1{]}

A 10-bit Charge-Redistribution ADC Consuming 1.9 uW at 1 MS/s {[}2{]}

A 6.3 uW 20 bit Incremental Zoom-ADC with 6 ppm INL and 1 uV Offset
{[}3{]}

A 12-Bit 1.25-GS/s DAC in 90 nm CMOS With \textgreater70 dB SFDR up to
500 MHz {[}4{]}

\phantomsection\label{refs}
\begin{CSLReferences}{0}{0}
\bibitem[\citeproctext]{ref-wu19}
\CSLLeftMargin{{[}1{]} }%
\CSLRightInline{W. Wu \emph{et al.}, {``A 28-nm 75-fsrms analog
fractional- \(N\) sampling PLL with a highly linear DTC incorporating
background DTC gain calibration and reference clock duty cycle
correction,''} \emph{IEEE Journal of Solid-State Circuits}, vol. 54, no.
5, pp. 1254--1265, 2019, doi:
\href{https://doi.org/10.1109/JSSC.2019.2899726}{10.1109/JSSC.2019.2899726}.
}

\bibitem[\citeproctext]{ref-elzakker10}
\CSLLeftMargin{{[}2{]} }%
\CSLRightInline{M. van Elzakker, E. van Tuijl, P. Geraedts, D. Schinkel,
E. A. M. Klumperink, and B. Nauta, {``A 10-bit charge-redistribution ADC
consuming 1.9 \(\mu\)w at 1 MS/s,''} \emph{IEEE Journal of Solid-State
Circuits}, vol. 45, no. 5, pp. 1007--1015, 2010, doi:
\href{https://doi.org/10.1109/JSSC.2010.2043893}{10.1109/JSSC.2010.2043893}.
}

\bibitem[\citeproctext]{ref-chae13}
\CSLLeftMargin{{[}3{]} }%
\CSLRightInline{Y. Chae, K. Souri, and K. A. A. Makinwa, {``A 6.3
\(\mu\)w 20 bit incremental zoom-ADC with 6 ppm INL and 1 \(\mu\)v
offset,''} \emph{IEEE Journal of Solid-State Circuits}, vol. 48, no. 12,
pp. 3019--3027, 2013, doi:
\href{https://doi.org/10.1109/JSSC.2013.2278737}{10.1109/JSSC.2013.2278737}.
}

\bibitem[\citeproctext]{ref-tseng11}
\CSLLeftMargin{{[}4{]} }%
\CSLRightInline{W.-H. Tseng, C.-W. Fan, and J.-T. Wu, {``A 12-bit
1.25-GS/s DAC in 90 nm CMOS with \(>\)70 dB SFDR up to 500 MHz,''}
\emph{IEEE Journal of Solid-State Circuits}, vol. 46, no. 12, pp.
2845--2856, 2011, doi:
\href{https://doi.org/10.1109/JSSC.2011.2164302}{10.1109/JSSC.2011.2164302}.
}

\end{CSLReferences}
